\section{Flat real POVMs: exact commuting rigidity and a descent obstruction}
For an unbiased pairwise-independent law on $n-1$ sign bits write
$v_x=(1,x)\in\{\pm1\}^n$ and keep only positive-weight atoms.
Then $\sum_xw_xv_xv_x^T=I_n$. Set
\[
 u_x=v_x/\sqrt n,\quad P_x=u_xu_x^T,\quad
 E_x=nw_xP_x=w_xv_xv_x^T,\quad\lambda_x=nw_x.
\]
The $E_x$ form a real rank-one positive operator-valued measurement
(POVM): they are positive semidefinite and sum to identity.
Here ``flat'' means that the unit vectors have entries of magnitude
$1/\sqrt n$, not merely that the effect weights are equal.
Distinct atoms with the first coordinate normalized positive give
distinct nonparallel directions. The POVM inequality $E_x\le I_n$
implies $0<\lambda_x\le1$.

\begin{theorem}[Exact projective and commuting criteria]
For such a distinct-direction flat POVM the following are equivalent:
the law is a uniform Hadamard law; every $E_x$ is a projection;
the effects commute pairwise; and
$\sum_{x,y}\|[E_x,E_y]\|_F^2=0$.
Its total projectivity defect is exactly
\[
 \sum_x\operatorname{tr}(E_x-E_x^2)=n-n^2C(w).
\]
Writing $s_{xy}=v_x\cdot v_y$ and $A=\sum_xE_x^2$, its ordered
commutator energy is
\[
 \mathcal J=2\sum_{x,y}w_x^2w_y^2s_{xy}^2(n^2-s_{xy}^2)
 =2\left[\operatorname{tr}(A^2)-
                 \sum_{x,y}(\operatorname{tr}(E_xE_y))^2\right].
\]
\end{theorem}
\begin{proof}
The eigenvalue of $E_x$ on its image is $\lambda_x$, so
$\operatorname{tr}E_x=nw_x$ and $\operatorname{tr}E_x^2=n^2w_x^2$.
Two nonzero rank-one symmetric effects commute precisely when their
directions are orthogonal or parallel: their products are scalar
multiples of $u_xu_y^T$ and $u_yu_x^T$. Parallelism is excluded here.
Thus commuting forces mutually orthogonal directions. Completeness
then requires exactly $n$ of them, each with eigenvalue one, giving
the normalized rows of a Hadamard matrix. Conversely those rows give
commuting projections. A sum of squared norms vanishes exactly when
each commutator does. Finally expand
$[E_x,E_y]=w_xw_ys_{xy}(v_xv_y^T-v_yv_x^T)$ and use
$\|v_x\|^2=\|v_y\|^2=n$ to obtain the energy formula.
\end{proof}
Aggregating duplicate directions is essential to the commuting
criterion: splitting one projector into several parallel effects
would otherwise give a commuting but nonprojective measurement.

\begin{theorem}[Uniform tight frames are uniformly almost commuting]
If the law is uniform on $m$ atoms, put $\lambda=n/m$. Then
\[
 \max_x\|E_x-E_x^2\|_{\rm op}=\lambda(1-\lambda),\qquad
 \max_{x,y}\|[E_x,E_y]\|_{\rm op}^2
 \le\lambda^3(1-\lambda),\qquad
 \mathcal J/n\le2\lambda^2(1-\lambda).
\]
For every positive integer $r$, the normalized projectivity defect
of the self-tensor POVM is $1-(nC(w))^r$.
No self-tensor power becomes projective or commuting unless the
original POVM was already a Hadamard measurement.
\end{theorem}
\begin{proof}
Write $c_{xy}=u_x\cdot u_y$. Tightness gives
$\lambda\sum_yc_{xy}^2=1$, and hence
$\sum_{y\ne x}c_{xy}^2=(1-\lambda)/\lambda$.
The rank-two commutator has operator norm
$\lambda^2|c_{xy}|\sqrt{1-c_{xy}^2}$ and squared Frobenius norm
$2\lambda^4c_{xy}^2(1-c_{xy}^2)$. Bounding $1-c_{xy}^2\le1$
and summing proves both estimates. Tensor collision is $C(w)^r$,
so the defect identity gives the displayed normalized value.
If some base commutator is nonzero, fix any other nonzero effect $E_z$.
The tensor effects with first factors $E_x,E_y$ and the remaining
factors $E_z$ have commutator
$[E_x,E_y]\otimes(E_z^2)^{\otimes(r-1)}\ne0$. This excludes
commuting activation, and projectivity is excluded by the defect.
\end{proof}

\subsection{Even flat dilations with vanishing overhead do not ensure descent}
The finite-dimensional real dilation is explicit. The $m\times n$
matrix $W$ with rows $\sqrt{w_x}v_x^T$ satisfies $W^TW=I_n$, and
$E_x=W^T|x\rangle\langle x|W$. Complete its orthonormal columns to
a real orthogonal matrix. This is the elementary finite-rank
Naimark construction; the added columns have no automatic flatness.

In fact the character construction established in the earlier
stability branch gives a stronger counterexample. Use $h=2^s$,
known Hadamard order $k<h$, $g=h-k\ge s$, and slopes
$T\subset\mathbb F_h$ of size $k$ containing zero, whose omitted
set spans $\mathbb F_h$ and has size $g<h/2$.
For seeds $(a,b)$, the observed blocks are all nonconstant characters
of $a+tb$, $t\in T$, plus the mixed characters $\chi_e(tb)$ for
$t\ne0$. The resulting flat real POVM has
$m=h^2$, $n=kh$, and $\lambda=k/h$.
All augmented columns are distinct Walsh characters on the two
seeds. Thus $W$ is literally $n$ columns of a normalized real Walsh
matrix of order $m$: an \emph{exact flat real dilation already exists}.
Its dimension overhead is $m/n-1=h/k-1$.

The source orthogonality graph is exactly $h$ disjoint cliques of
size $h$, and completing these cliques using a known order-$k$
Hadamard matrix gives
\[
 \inf_{q\in\mathcal H_n}\|w-q\|_{\rm TV}=1-1/k,
\]
where $\mathcal H_n$ denotes mixtures of normalized real Hadamard
laws. Along the established known-order subsequences $k/h\to1$ and
$g\to\infty$, both effectwise defects, the normalized commutator
energy and the relative dilation overhead tend to zero, while this
distance tends to one. Therefore there is no dimension-independent
flat-real projectivization estimate tending to zero from these
defects. The counterexample even starts with an explicit flat
Naimark dilation, so flat dilation alone cannot supply the descent.

The distance here uses the common outcome labels given by actual
normalized sign directions. On the maximally mixed state $I_n/n$,
the measurement outcome probabilities are exactly $w_x$. Thus the
same lower bound applies to any measurement distance that dominates
the outcome TV distance for that input with these labels. No claim
is made for arbitrary relabeling, postprocessing, or unrestricted
complex projective measurements. This distinction matters: the
dimension-deficient Naimark and POVM-simulation literature, for
example Kotowski--Oszmaniec, allows different projective directions
and noise/postprocessing; see
\url{https://arxiv.org/abs/2501.09339}. Their usage of nearly flat
effects concerns weights, not the entrywise flatness required here.
The present identities and classical dilation are not claimed new;
the quantitative obstruction inherits the earlier sign-law theorem
and does not prove missing Hadamard existence.
